\section{Step 11: Prove the exact payload extent}
Read the four-byte payload length at header offset eight by the same
big-endian method. The admitted total count includes the header. Subtract sixteen
from that count and require equality with the declared payload length. This
rejects both missing bytes and an unexplained suffix. The outer count bound
already limits the payload to 28\,672 bytes.

Fill \texttt{record\_extent}. Save the payload length at stack offset 96. Set
\texttt{\$s5} to source plus sixteen and \texttt{\$s6} to that address plus the
length. The latter is an exclusive end address: it names the boundary after the
last payload byte. Later checks compare a remaining byte count before reading;
they must not load from the exclusive end itself.

\paragraph{Observe} Prepare \texttt{s11-1} and stop before
\texttt{record\_extent\_ready}. At \texttt{record\_extent\_ready}, the payload length in
\texttt{\$t0} is eleven. \texttt{\$s5} and \texttt{\$s6} are source plus
16 and source plus 27. The source is the address passed in \texttt{\$a0},
three bytes after \texttt{record\_data}.

\Needspace{7\baselineskip}
\paragraph{Check} What does the payload end pointer name?
\begin{itemize}
  \item[A.] The last payload byte.
  \item[B.] The checksum byte.
  \item[C.] The first address beyond the payload.
\end{itemize}
